\providecommand{\Res}[2]{\qty[round-mode=figures,round-precision=3]{\fpeval{#1}}{#2}}
\providecommand{\Val}[1]{\num[round-mode=figures,round-precision=4]{\fpeval{#1}}}
\providecommand{\ResSci}[2]{\qty[round-mode=figures,round-precision=3,exponent-mode=scientific]{\fpeval{#1}}{#2}}
\subsubsection*{Question a) Angle du premier minimum}
\begin{enumerate}[label=\arabic*)]
    \item \textbf{Variables connues et recherchées :}
    \begin{itemize}[label= -]
        \item $\lambda = \qty[exponent-mode=threshold]{\vsi{lam}}{m}$, $a = \qty[exponent-mode=threshold]{\vsi{larg}}{m}$, $m = 1$
        \item Variable recherchée : $\theta_1$
    \end{itemize}
    \item \textbf{Équations pertinentes :} minimums de diffraction d'une fente :
    \[ a\sin\theta = m\lambda \]
    \item \textbf{Isolation de la variable recherchée :}
    \[ \theta_1 = \arcsin\left(\frac{\lambda}{a}\right) \]
    \item \textbf{Résolution numérique :}
    \[ \theta_1 = \arcsin\left(\frac{\num[exponent-mode=threshold]{\vsi{lam}}}{\num[exponent-mode=threshold]{\vsi{larg}}}\right) = \Res{asind(\vsi{lam}/\vsi{larg})}{\degree} \]
    \item \textbf{Réponse :} le premier minimum est à \textbf{\Res{asind(\vsi{lam}/\vsi{larg})}{\degree}}.
\end{enumerate}

\subsubsection*{Question b) Largeur du maximum central}
\begin{enumerate}[label=\arabic*)]
    \item \textbf{Variables connues et recherchées :} $L = \qty{\v{L}}{m}$, $\theta_1$ ; variable recherchée : $W$
    \item \textbf{Équations pertinentes :} le maximum central s'étend du premier minimum d'un côté au premier minimum de l'autre : $y_1 = L\tan\theta_1$
    \item \textbf{Isolation de la variable recherchée :}
    \[ W = 2y_1 = 2L\tan\theta_1 \]
    \item \textbf{Résolution numérique :}
    \[ W = 2 \cdot \num{\v{L}} \cdot \tan(\Val{asind(\vsi{lam}/\vsi{larg})}^\circ) = \Res{2000*\v{L}*tand(asind(\vsi{lam}/\vsi{larg}))}{mm} \]
    \item \textbf{Réponse :} le maximum central mesure \textbf{\Res{2000*\v{L}*tand(asind(\vsi{lam}/\vsi{larg}))}{mm}} de large.
\end{enumerate}

\subsubsection*{Question c) Largeur d'un maximum secondaire}
Pour de petits angles, $y_m \approx m\lambda L/a$, donc deux minimums voisins sont séparés de
\[ \Delta y = \frac{\lambda L}{a} = \frac{\num[exponent-mode=threshold]{\vsi{lam}} \cdot \num{\v{L}}}{\num[exponent-mode=threshold]{\vsi{larg}}} = \Res{1000*\vsi{lam}*\v{L}/\vsi{larg}}{mm} \]
Un maximum secondaire mesure \textbf{\Res{1000*\vsi{lam}*\v{L}/\vsi{larg}}{mm}}, soit la moitié de la largeur du maximum central.
